\chapter{Linear Stability and Frequency-Domain Methods}
\label{ch:linear}

Stability is, in the first instance, a question about \emph{small} perturbations:
does a tiny disturbance grow or decay? This is the domain of linear stability
analysis. We linearise the coupled model, introduce the zero-power and
feedback transfer functions, and apply the Nyquist criterion, which handles the
transport delays of real reactors more gracefully than eigenvalues alone.

\section{Linearisation about an operating point}

Let the reactor sit at steady power $p_0$ with steady temperatures, and perturb:
$p=p_0+\delta p$, $\zeta_i=\zeta_{i0}+\delta\zeta_i$, $\Tf=\Tf^0+\delta\Tf$, and
so on. Substituting into the coupled equations of Chapter~\ref{ch:coupled} and
discarding products of small quantities gives a linear system
\begin{equation}
\frac{\dd}{\dd t}\,\delta\mathbf{x} = \mathbf{A}\,\delta\mathbf{x}
   + \mathbf{b}\,\delta\rho_{ext},
\end{equation}
where $\delta\mathbf{x}=(\delta p,\delta\zeta_1,\dots,\delta\zeta_6,
\delta\Tf,\delta\Tc)^{\!\top}$. The crucial linearised coupling is in the power
equation: the term $(\rho-\beta)p/\Lambda$ contributes
$p_0\,\delta\rho/\Lambda$, and $\delta\rho = \alpha_D\,\delta\Tf +
\alpha_c\,\delta\Tc$. Thus the feedback coefficients enter $\mathbf A$ multiplied
by $p_0/\Lambda$: \emph{feedback is more potent at higher power}, an observation
with direct bearing on why the RBMK was most dangerous at low power.

The operating point is asymptotically stable if and only if every eigenvalue of
$\mathbf A$ has negative real part. For the bare point-kinetics block (no
feedback) the eigenvalues are exactly the inhour roots of
Chapter~\ref{ch:pk}; adding negative feedback pushes the dominant root further
into the left half-plane (more stable), while positive feedback pulls it toward,
and possibly across, the imaginary axis (instability).

\section{The zero-power transfer function}

Frequency-domain analysis is more illuminating for systems with delays. Laplace
transform the point kinetics about $p_0$ with input $\delta\rho$ and output
$\delta p/p_0$; eliminating the precursors gives the \textbf{zero-power reactor
transfer function}
\begin{equation}
Z(s) = \frac{\delta p/p_0}{\delta\rho}
     = \frac{1}{\,\Lambda s + \displaystyle\sum_{i}\frac{\beta_i s}{s+\lambda_i}\,}.
\label{eq:zpr}
\end{equation}
Its denominator is exactly the inhour function \eqref{eq:inhour}; its poles are
the inhour roots. The name ``zero-power'' signals that no thermal feedback is
included: it is the open-loop response of the neutronics alone. At low frequency
$Z(s)$ behaves like an integrator with a long effective time constant set by the
precursors; at high frequency $Z(s)\to 1/(\Lambda s)$, the fast prompt response.
The transition between these regimes, around the delayed-neutron frequencies, is
where the reactor's dynamic character lives.

\section{Closing the loop: the feedback transfer function}

Let the feedback path have transfer function $H(s)$ relating a power perturbation
to the feedback reactivity it produces, $\delta\rho_{fb}=-H(s)\,\delta p/p_0$
(the minus sign for negative feedback). For a single lumped fuel node with
coefficient $\alpha_D<0$ and thermal time constant $\tau_f$,
\begin{equation}
H(s) = \frac{|\alpha_D|\,K}{1+\tau_f s},
\end{equation}
a first-order lag: fast feedback at low frequency, attenuated and phase-lagged at
high frequency. The closed-loop transfer function from external reactivity to
power is
\begin{equation}
G(s) = \frac{Z(s)}{1 + Z(s)H(s)},
\end{equation}
and stability is governed by the roots of the characteristic equation
$1+Z(s)H(s)=0$. The product $Z(s)H(s)$ is the \emph{open-loop gain}; everything
hinges on its behaviour.

\section{The Nyquist criterion}

The Nyquist criterion reads stability directly off a plot of the open-loop gain
$Z(j\omega)H(j\omega)$ as $\omega$ runs over all frequencies. The closed loop is
stable provided the Nyquist locus does not encircle the critical point $-1$ in the
complex plane. Physically, encirclement of $-1$ means there is a frequency at
which the feedback, after its phase lag, returns \emph{in phase} with the
disturbance and with gain exceeding unity---positive feedback at that frequency,
hence growth.

This is where \emph{phase lag} becomes dangerous even for nominally negative
feedback. If the feedback is negative at zero frequency but accumulates enough
phase lag (from thermal transport delays, coolant transit, etc.) that at some
frequency it is delayed by half a cycle, it acts as \emph{positive} feedback at
that frequency. If the loop gain there exceeds one, the reactor oscillates with
growing amplitude. Boiling-water reactors, with their void feedback transported
by the coolant, are the classic example: they can exhibit lightly damped or even
growing power--flow oscillations under certain low-flow, high-power conditions,
and operating limits are drawn specifically to keep the decay ratio safely below
one.

\begin{definition}[Decay ratio]
The decay ratio is the ratio of the amplitudes of two successive peaks of an
oscillatory response to a perturbation. A decay ratio below $1$ indicates a
damped, stable response; equal to $1$ marks the stability boundary; above $1$ the
oscillation grows. BWR stability is monitored and licensed in these terms,
typically with a limit well below unity.
\end{definition}

\section{Gain and phase margins}

The robustness of stability is quantified by the \emph{gain margin} (how much the
loop gain could increase before the locus reaches $-1$) and the \emph{phase
margin} (how much additional phase lag could be tolerated). A reactor designed for
inherent stability carries comfortable margins so that variations over the fuel
cycle, across power levels, and under off-normal conditions never erode stability.
The RBMK at low power had, in effect, negative margins in the relevant
configuration: its dominant feedback was positive, so no amount of phase
bookkeeping could save it---the locus enclosed $-1$ outright.

\section{Power dependence and the low-power trap}

Because feedback enters the loop gain multiplied by $p_0$, a reactor's stability
character changes with power. A \emph{positive}-feedback reactor like the RBMK is
most unstable at \emph{low} power: there the steam void fraction is small and most
sensitive, so the void coefficient is largest, while the prompt Doppler feedback
(which scales with the fuel-temperature changes that low power barely produces) is
weakest. The result is a window of operation---low power, high positive void
worth, depleted control margin---in which the machine is acutely unstable. The
Chernobyl test drove the reactor straight into that window.

\keypoint{Linear stability is decided by the eigenvalues of the linearised
system, or equivalently by whether the Nyquist locus of the open-loop gain
$Z(s)H(s)$ encircles $-1$. Negative feedback with small phase lag keeps the locus
clear; positive feedback, or negative feedback with excessive transport lag, can
enclose $-1$ and produce growing oscillations or divergence. Stability margins
must hold across all power levels and cycle conditions---a requirement the RBMK
violated at low power.}



\section{State-space form and the Routh--Hurwitz criterion}
\label{sec:statespace}

To make ``stability'' fully explicit, reduce to the one-precursor,
one-temperature model that retains the essential feedback physics:
\begin{equation}
\dot n=\frac{\rho-\beta}{\Lambda}n+\lambda C,\quad
\dot C=\frac{\beta}{\Lambda}n-\lambda C,\quad
\dot T=\frac{n}{\mu}-\theta T,\quad \rho=\rho_0+\alpha T,
\end{equation}
with heat capacity $\mu>0$, cooling rate $\theta>0$, and feedback coefficient
$\alpha$. Linearising about the critical equilibrium $(n_0,C_0,T_0)$ with
$\rho_0=0$ gives the Jacobian
\begin{equation}
\mathbf A=
\begin{pmatrix}
-\dfrac{\beta}{\Lambda} & \lambda & \dfrac{\alpha n_0}{\Lambda}\\[5pt]
\dfrac{\beta}{\Lambda} & -\lambda & 0\\[5pt]
\dfrac{1}{\mu} & 0 & -\theta
\end{pmatrix},
\end{equation}
whose characteristic polynomial is $\det(s\mathbf I-\mathbf A)=s^3+a_1s^2+a_2s+a_3$
with
\begin{equation}
a_1=\lambda+\theta+\frac{\beta}{\Lambda},\quad
a_2=\theta\Big(\lambda+\frac{\beta}{\Lambda}\Big)-\frac{\alpha n_0}{\Lambda\mu},\quad
a_3=-\frac{\alpha n_0\lambda}{\Lambda\mu}.
\label{eq:rh-coeffs}
\end{equation}

\begin{theorem}[Routh--Hurwitz, cubic]
\label{thm:rh}
The polynomial $s^3+a_1s^2+a_2s+a_3$ (real coefficients) has all roots in the open
left half-plane if and only if $a_1>0$, $a_3>0$, and $a_1a_2>a_3$.
\end{theorem}

\begin{proposition}[Feedback sign decides stability]
\label{prop:stab}
With $\lambda,\theta,\beta,\Lambda,\mu,n_0>0$, the equilibrium is asymptotically
stable if and only if $\alpha<0$.
\end{proposition}
\begin{proof}
By \eqref{eq:rh-coeffs}, $a_1>0$ always. Now $a_3=-\alpha n_0\lambda/(\Lambda\mu)$,
so $a_3>0\iff\alpha<0$. If $\alpha<0$ then also $a_2=\theta(\lambda+\beta/\Lambda)
+|\alpha|n_0/(\Lambda\mu)>0$, and
\[
a_1a_2-a_3
=\underbrace{(\lambda+\theta+\tfrac{\beta}{\Lambda})\,\theta(\lambda+\tfrac{\beta}{\Lambda})}_{>0}
+\frac{|\alpha|n_0}{\Lambda\mu}\Big[(\lambda+\theta+\tfrac{\beta}{\Lambda})-\lambda\Big]
=(\cdots)+\frac{|\alpha|n_0}{\Lambda\mu}\Big(\theta+\tfrac{\beta}{\Lambda}\Big)>0,
\]
so all three Routh--Hurwitz conditions of Theorem~\ref{thm:rh} hold and the
equilibrium is asymptotically stable. Conversely, if $\alpha>0$ then $a_3<0$,
violating the necessary condition $a_3>0$, so at least one root lies in the closed
right half-plane and the equilibrium is unstable. ($\alpha=0$ is the neutral case,
$a_3=0$, with a root at the origin.)
\end{proof}

Proposition~\ref{prop:stab} is the precise version of the book's thesis: with all
thermal and kinetic parameters positive, \emph{the sign of the feedback
coefficient alone} decides linear stability. The RBMK's low-power state had the
effective $\alpha>0$ and was therefore linearly unstable.

\paragraph{Lyapunov viewpoint.}
By Theorem~\ref{thm:lyap} (Appendix~\ref{app:math}), $\alpha<0$ guarantees a
quadratic Lyapunov function $V(\mathbf x)=\mathbf x^{\!\top}\mathbf P\mathbf x$ with
$\mathbf P\succ0$ solving $\mathbf A^{\!\top}\mathbf P+\mathbf P\mathbf A=-\mathbf
Q$ for any $\mathbf Q\succ0$; then $\dot V=-\mathbf x^{\!\top}\mathbf Q\mathbf x<0$
certifies decay of every small perturbation. Stability is thus not merely
spectral but energy-dissipative.

\section{The Nyquist criterion, precisely}
\label{sec:nyquist-precise}

For the full delayed system the loop gain $L(s)=Z(s)H(s)$ is meromorphic, and
stability is read off by the argument principle.

\begin{theorem}[Nyquist]
\label{thm:nyquist}
Let $L(s)$ be the open-loop transfer function with $P$ poles in the open right
half-plane. As $s$ traverses the Nyquist contour (the imaginary axis closed by an
infinite right semicircle), let $N$ be the number of clockwise encirclements of
the point $-1$ by $L(s)$. Then the number of closed-loop poles in the right
half-plane is $Z=N+P$. The closed loop is stable iff $Z=0$, i.e.\ iff
$N=-P$.
\end{theorem}
\begin{proof}[Sketch]
Closed-loop poles are zeros of $1+L(s)$. By the argument principle, as $s$ runs
over the closed contour, the net change in $\arg(1+L(s))$ equals $2\pi(Z-P)$,
where $Z,P$ count zeros and poles of $1+L$ (equivalently poles of $L$) inside the
contour. The change in $\arg(1+L)$ is $2\pi$ times the clockwise encirclements of
the origin by $1+L$, i.e.\ of $-1$ by $L$; hence $N=Z-P$.
\end{proof}

For an open-loop-stable reactor ($P=0$) the criterion reduces to the familiar
statement: stability iff the locus $L(j\omega)$ does not encircle $-1$. Phase lag
in $H(s)$ from thermal transport can rotate the locus until it encircles $-1$,
which is the frequency-domain face of Proposition~\ref{prop:stab}'s sign change.


\section*{Exercises}
\addcontentsline{toc}{section}{Exercises}
\begin{enumerate}[leftmargin=*,itemsep=2pt]
\item Write the zero-power transfer function \eqref{eq:zpr} and identify its poles.
\item State the Nyquist stability criterion and interpret encirclement of $-1$ physically.
\item Explain how phase lag can turn nominally negative feedback into destabilising feedback at some frequency.
\item Define the decay ratio and state the stability condition in terms of it.
\item Explain why a positive-feedback reactor is most unstable at low power.
\end{enumerate}
